Die vorliegende Arbeit erfolgt in Kooperation mit der snoopmedia GmbH und nutzt Ergebnisse sowie Infrastrukturen aus einem vorangegangenen Forschungsprojekt. Im Rahmen dieser 
Zusammenarbeit steht eine etablierte Crawling-Pipeline sowie ein umfangreicher Datensatz an Rohdaten aus dem Bereich der Lebensmittel-Produktinformationen zur Verfügung. Die Arbeit 
ist dabei in eine bestehende Systemlandschaft eingebettet, wodurch die Anforderungen an die Struktur und Qualität der Zielformate bereits fest definiert sind. Die Evaluation des zu
 entwickelnden Lösungsansatzes erfolgt auf Basis eines repräsentativen Datenausschnitts aus diesem Bestand.